\documentclass[11pt,letterpaper]{article}
\usepackage{technical-user-manual}
\setDocTitle{Chapter-to-Operations Crosswalk}
\setDocSubtitle{Source foundations to modern practices}
\setDocVersion{1.0}
\setDocStatus{Technical education / proposed reference}
\setDocOwner{CI/CD Engineering}
\setDocDate{30 October 2026}
\setDocSummary{Supplementary engineering practice}
\begin{document}
\makedoctitle
\makedocsummary
\tableofcontents
\clearpage
\section{Purpose}
Book chapters map to operations and evidence without claiming that the 2007 source includes current GitHub tooling.

\section{Chapter 1: Getting Started}
Opening printed page 3. Core operational transfer: commit-to-build pipeline.

Laboratory demonstration: Construct the minimal GitHub Actions pipeline in laboratory 01 and record the first green commit.

Crosswalk gate: attach repository revision, tool/version, actual run result and the responsible reviewer.

\section{Chapter 2: Introducing Continuous Integration}
Opening printed page 23. Core operational transfer: short integration cadence and repair policy.

Laboratory demonstration: Execute the same test suite locally and as a PR check; document divergence and resolution.

Crosswalk gate: attach repository revision, tool/version, actual run result and the responsible reviewer.

\section{Chapter 3: Reducing Risks Using CI}
Opening printed page 47. Core operational transfer: risk-based gate mapping.

Laboratory demonstration: Inject one failing unit test and one missing configuration value and compare time-to-detection.

Crosswalk gate: attach repository revision, tool/version, actual run result and the responsible reviewer.

\section{Chapter 4: Building Software at Every Change}
Opening printed page 65. Core operational transfer: single-command staged builds and timing.

Laboratory demonstration: Collect ten pipeline durations before and after a cache or test-partition change; record confounders.

Crosswalk gate: attach repository revision, tool/version, actual run result and the responsible reviewer.

\section{Chapter 5: Continuous Database Integration}
Opening printed page 107. Core operational transfer: ephemeral database migration.

Laboratory demonstration: Run forward migrations in an ephemeral SQLite sandbox and verify test fixtures and expected columns.

Crosswalk gate: attach repository revision, tool/version, actual run result and the responsible reviewer.

\section{Chapter 6: Continuous Testing}
Opening printed page 129. Core operational transfer: test pyramid and regression evidence.

Laboratory demonstration: Implement and execute the Python, Java and C++ sample test suites and compare their signals.

Crosswalk gate: attach repository revision, tool/version, actual run result and the responsible reviewer.

\section{Chapter 7: Continuous Inspection}
Opening printed page 161. Core operational transfer: inspection triage and coverage interpretation.

Laboratory demonstration: Introduce a complexity hotspot or unsafe deserialization sample in an isolated branch and triage without deploying.

Crosswalk gate: attach repository revision, tool/version, actual run result and the responsible reviewer.

\section{Chapter 8: Continuous Deployment}
Opening printed page 189. Core operational transfer: artifact promotion and rollback.

Laboratory demonstration: Complete the deployment architecture specification and demonstrate a simulated failed release decision.

Crosswalk gate: attach repository revision, tool/version, actual run result and the responsible reviewer.

\section{Chapter 9: Continuous Feedback}
Opening printed page 203. Core operational transfer: notification routing and SLO reporting.

Laboratory demonstration: Generate a failure notification from an example run and measure acknowledgement and restoration times.

Crosswalk gate: attach repository revision, tool/version, actual run result and the responsible reviewer.

\section{Supplementary standards crosswalk}
GitHub Actions workflow policy maps to M07--M08; supply chain and SLSA v1.2 to M06--M08; measurement definitions and maturation assessment to M09--M10. These mappings are the author's engineering extensions, not normative statements from the book.

\end{document}
